\documentclass[10pt]{article}

\usepackage[utf8]{inputenc}

\usepackage[T1]{fontenc}

\usepackage{amsmath}

\usepackage{amsfonts}

\usepackage{amssymb}

\usepackage[version=4]{mhchem}

\usepackage{stmaryrd}



\begin{document}

зал первую теорему теории простых чисел - бесконечность числа простых чисел. Несколько позднее Эратосфеном был найден метод получения простых чисел, к-рый стал наз. Эратосбена решетом.



Систематизация проблем Ч. т. и методов их решений проведена Диофантом (3 в.) в его «Арифметике», где, в частности, дано решение в рациональных числах многих алгебраич. уравнений 1-й и 2-й степени с целыми коәффициентами от нескольких неизвестных.



В Китае, начиная со 2 в., в связи с календарными расчетами, возникла задача определения наименьшего целого числа, дающего при делении на заданные тисла заданные остатки, к-рая была решена Сунь-цзы (26 вв.) и Цинь Цзютао (13 в.) .



В Индии Брамагупта (7 в.) и Бхаскара (12 в.) дали общие методы решения в целых числах неопределенных уравнений 1-й степени с двумя неизвестными и уравнений вида $a x^{2}+b=c y^{2}, x y=a x+b y+c$.



Расцвет Ч. ч. начался в Европе с работ П. Ферма (P. Fermat, 17 в.). Он исследовал решения многих уравнений в целых числах, в частности высказал гипотезу, что уравнение $x^{n}+y^{n}=z^{n}, n>2$, не имеет решентй в натуральных числах $x, y, z$ (Ферма теорема великая), доказал, что простые числа вида $4 n+1$ являются суммою двух квадратов, высказал одно из основных утверждений теории сравнений: $a p-a$ делится на $p$, если р-простое число (Ферма малая теорема).



Исключительно важный вклад в Ч. т. внес Л. Эйлер (L. Euler). Он доказал теорему Ферма при $n=3$, малую теорему Ферма и ее обобщение, целый ряд теорем о представлении чисел бинарными квадратичными бормами. Л. Эйлер был первым, кто для решения задач Ч. т. привлек средства математич. анализа, что привело к созданию аналитической теории чисел. Производящие функции Эйлера явились источником кругового метода Харди - Литлвуда - Рамануджана и триигоометрических сумм метода И. М. Виноградова- основных методов современной аддитивной теории чисел, а функция $\zeta(s)$, введенная Л. Эйлером, и ее обобщения составляют основу современных аналитич. методов исследования проблем распределения простых чисел.



В переписке Х. Гольдбаха (Ch. Goldbach) с Л. Эйлером были поставлены три знаменитые проблеме: всякое нечетное $N \geqslant 5$ есть сумма трех простых, четное $N \geqslant$ $\geqslant 4$ есть сумма двух простых и нечетное $N=p+2 n^{2}$, где $n$ - целое, $p$ - простое. Первая проблема решена И. М. Виноградовым (1937), две другие проблемы не решены (1984).



K. Гаусc (C. Gauss) создал основные методы и заверпил построение теории сравнений, доказал взаимности закон квадратичных вычетов, сформулированный Л. Эйлером, заложил основы теории представления чисел квадратичными форормами $a x^{2}+b x y+c y^{2}$ и формами высших степеней со многими переменными, ввел т. н. Гаycca сумми



\[

S=\sum_{n=0}^{m-1} e^{2 \pi i \frac{a n^{2}}{m}},

\]



к-рые явились первыми тригонометрич. суммами, и показал их полезность в решении задач Ч. т. Если до $К$. Гаусса Ч. т. представляла собой собрание отдельных результатов и идей, то после его работ она начала развиваться в различных направлениях как стройная теория.



Большой вклад в развитие Ч. т. внесли многие ученые $19-20$ вв.



Особенностью и привлекательностью Ч. т. является простота и доступность формулировок большинства проблем и трудность их решений. Напр., проблема бесконечности числа простых чисел близнецов была поставлена еще Евклидом и не решена (1984). Отдельные проблемы Ч. т. стали источником больших самостоятельных направлениї математики. Такими являются: теория простых чисел и связанная с ней теория дзетабункций и Дирихле рядов, теория диобантовых уравнений, аддитивная теория чисел, метрическая теория чисел, теория алгебраических и трансцендентных чисел, алгебраическая теория чисел, теория диофантовых приближений, чисел теория вероятностная, геометрия чисел. Напр., источником аналитич. Ч. т. послужили проблема распределения простых чисел в натуральном ряде и проблема представления натуральных чисел суммою слагаемых определенного вида. Решение уравнений в целых числах, к-рые стали наз. диофантовыми, в частности великая теорема Ферма, послужили источником алгебраич Ч. т. Проблема построения с помощы циркуля и линейки круга единичной площади (квадратура круга) привела к вопросу об арифметич. природе числа $\pi$, а вместе с ним - к созданию теории алгебраических и трансцендентных чисел.



Все названные направления Ч. т. переплетаются между собой, дополняя и обогашая друг друга.



Jum.: [1] в и н о г p а д о в И. М., Основы теории чисел, 9 изд., М., 1981; [2] е г о ж е, Метод тригонометрических сумм, метрических сумм, М., 1976; [4] К а р а ч у 6 a А. А., Основы аналитической теории чисел, 2 изд., М., 1983: [5] Б о р ев и ч З. И., ІІ а ф а р е в и ч И. Р., Теория чисел, 2 изд., М.,

$1972 ;$ [6] Д в в е н п о р т Г., Мультипликативная теория чисел, пер. с англ., М., 1971; ' [7] Ч а н д p а с е к х а р a н K., Введение в аналитическую теорию чисел, пер. с англ., М., 1974; [9] Д и р и х л е П. Г. Л., Јекции по теории чисел, пер. с нем., M. - Л., 1936; [10] Т и т ч м a p III E. К., Теория дзета-функции Римана, пер. с англ., М., 1953; [11] В е н к о в Б. А., Элементарная теория чисел, М.- Л., 1937.



А. А. Карачуба и материалы одночменной статьи БСЭ-2. ЧИСЕЛ ТЕОРИЯ в е р о я т о о с т н а я - в широком смысле раздел теории чисел, в к-ром используются идеи и методы теории вероятностей.



Под вероятностной Ч. т. в узком смысле понимается статистич. теория распределения значений арифметических функций.



Подавляющее большинство арифметич. функций, изучаемых в теории чисел, являются аддитивными или мультипликативными. Их значения обычно распределены очень сложно. Если проследить за изменением значений таких функций, когда аргумент пробегает последовательные натуральные числа, получится весьма хаотическая картина, к-рая обычно наблюдается при рассмотрении аддитивных свойств целых чисел совместно с мультипликативными. В классич. исследованиях при рассмотрении распределения значений действительных арифметич. функций $f(m)$ обычно изучалось асимптотич. поведение самой функции $f(m)$ или ее среднего значения. В первом случае ищутся простые функции $\psi_{1}(m), \psi_{2}(m)$, чтобы было $\psi_{1}(m) \leqslant f(m) \leqslant \psi_{2}(m)$ для всех $m$ или хотя бы для всех достатотно больших $m$. Напр., если $\omega(m)$ означает число всех различных простых делителей числа $m$, то $\omega(m) \geqslant 1$ для всех $m>1$, $\omega(m) \leqslant 2(\ln \ln m)^{-1} \ln m$ при $m \geqslant m_{0} ;$



\[

\begin{gathered}

\lim _{m \rightarrow \infty} \inf \omega(m)=1, \\

\lim _{m \rightarrow \infty} \sup \omega(m)(\ln m)^{-1} \ln \ln m=1 .

\end{gathered}

\]



Во втором случае рассматривается поведение



\[

\frac{1}{n} \sum_{m=1}^{n} f(m)

\]



Для $\omega(m)$ среднее значение (1) равно $(1+o(1) \ln \ln n)$. Решение как первой, так и второй задачи в общем случае дает мало информации о поведении функции $f(m)$, об ее колебаниях. Функция может значительно отклоняться от своего среднего значения. При этом оказывается, что большие отклонения встречаются вообще довольно редко. Ставится задача отыскания границ, в к-рых могут колебаться значения функции $f(m)$ для подавляющего большинства значений аргумента.





\end{document}